\subsection{White Streaks}

\subsubsection{Описание решения}

В этой задаче требуется найти суммарное количество горизонтальных и вертикальных белых полос в произвольном прямоугольнике. Если её решать «в лоб» обходом каждой клетки прямоугольника, то мы не уложимся по времени выполнения алгоритма. Следовательно, нужно найти оптимизацию.

Заметим, что на ответ влияют лишь чёрные клетки. То есть количество белых клеткок между двумя чёрными однозначно определено. Сфокусируемся на их обработке.

Отсортируем все чёрные клетки сначала по абсциссе, затем по ординате. При последовательном обходе такого массива чёрные клетки одной колонки будут соседними. Следовательно, количество промежутков между такими клетками равно количеству вертикальных полос в данной колонке. Пройдёмся по каждой колонке прямоугольника и посчитаем их.

Стоит заметить, что в таком случае краевые ситуации не учитываются \textit{(к примеру, промежуток между верхней стороной прямоугольника и первой чёрной клеткой данной колонки, если он есть, не считается)}. Для разрешения этой проблемы добавим в массив чёрных клеток фиктивные элементы, которые расположены по краям данного прямоугольника. Они не увеличат число белых полос, но позволят сохранить логику алгоритма.

Аналогично поступим с подсчётом горизонтальных белых полос: отсортируем чёрные клетки сначала по ординате, затем по абсциссе, и посчитаем число горизонтальных промежутков.

В описанном алгоритме есть недостаток: одиночные белые клетки, которые окружены со всех сторон чёрными клетками, учитываются дважды \textit{(как горизонтальная и вертикальная полоса)}. Это недопустимо, поэтому на этапе подсчёта вертикальных полос будем вести множество клеток, которые подозрительны на «одиночество»: вертикальное расстояние между соседними чёрными клетками должно равняться двум. Если на этапе подсчёта горизонтальных полос гипотеза подтвердилась, то учитываем клетку лишь один раз. Иначе это обычная горизонтальная полоса, которая ограничена сверху и снизу произвольными чёрными клетками.

\subsubsection{Асимптотическая сложность}

\textit{Временная сложность решения} --- линейно-логарифмическая \( \mathcal{O}((N+M+K)\cdot\log(N+M+K)) \). В алгоритме дважды используется сортировка линейно-логарифмической сложности массива длины \( 2\cdot(N+M+2)+K \).

После каждой сортировки происходит линейный обход полученного массива. На каждой итерации выполняются операции не более чем логарифмической сложности \textit{(в связи с использованием красно-чёрного дерева для учёта подозрительных клеток)}.

\textit{Пространственная сложность решения} --- линейная \( \mathcal{O}(N+M+K) \). Массив чёрных клеток материализован, к нему добавлены в явном виде граничные клетки данного прямоугольника. Промежуточные вычисления занимают константное количество памяти.

\subsubsection{Листинг кода}

\begin{lstlisting}
#include <algorithm>
#include <cstddef>
#include <iostream>
#include <set>
#include <vector>

namespace {
struct Point {
  std::size_t x;
  std::size_t y;

  bool operator<(const Point& other) const {
    if (x != other.x) {
      return x < other.x;
    }
    return y < other.y;
  }
};

void AddBounds(std::vector<Point>& arr, const std::size_t x_bound, const std::size_t y_bound) {
  for (std::size_t y_idx = 0; y_idx <= y_bound + 1; y_idx++) {
    for (const std::size_t x_idx : {static_cast<std::size_t>(0), x_bound + 1}) {
      arr.push_back({x_idx, y_idx});
    }
  }
  for (std::size_t x_idx = 1; x_idx < x_bound + 1; x_idx++) {
    for (const std::size_t y_idx : {static_cast<std::size_t>(0), y_bound + 1}) {
      arr.push_back({x_idx, y_idx});
    }
  }
}

bool SortByY(const Point& lpt, const Point& rpt) {
  if (lpt.y != rpt.y) {
    return lpt.y < rpt.y;
  }
  return lpt.x < rpt.x;
}
}  // namespace

int main() {
  std::size_t m_count = 0;
  std::size_t n_count = 0;
  std::size_t k_count = 0;
  std::cin >> m_count >> n_count >> k_count;

  std::vector<Point> black_points(k_count);
  black_points.reserve(k_count + 2 * (n_count + m_count + 2));
  for (std::size_t i = 0; i < k_count; i++) {
    std::cin >> black_points[i].x >> black_points[i].y;
  }
  AddBounds(black_points, m_count, n_count);

  std::size_t white_streaks = 0;
  std::set<Point> suspicious_points = {};
  std::sort(black_points.begin(), black_points.end());
  for (std::size_t i = 0; i < black_points.size() - 1; i++) {
    const Point& lpt = black_points[i];
    const Point& rpt = black_points[i + 1];
    if (lpt.x != rpt.x) {
      continue;
    }
    const auto diff = static_cast<std::ptrdiff_t>(rpt.y - lpt.y);
    if (diff == 1) {
      continue;
    }
    if (diff == 2) {
      suspicious_points.insert({lpt.x, lpt.y + 1});
      continue;
    }
    white_streaks++;
  }

  std::sort(black_points.begin(), black_points.end(), SortByY);
  for (std::size_t i = 0; i < black_points.size() - 1; i++) {
    const Point& lpt = black_points[i];
    const Point& rpt = black_points[i + 1];
    if (lpt.y != rpt.y) {
      continue;
    }
    const auto diff = static_cast<std::ptrdiff_t>(rpt.x - lpt.x);
    if (diff == 1) {
      continue;
    }
    if (diff == 2) {
      const Point sus_point = {lpt.x + 1, lpt.y};
      if (suspicious_points.count(sus_point) != 0U) {
        suspicious_points.erase(sus_point);
        white_streaks++;
      }
      continue;
    }
    white_streaks++;
  }
  std::cout << white_streaks;
  return 0;
}
\end{lstlisting}